\section{Nutzung von Künstlicher Intelligenz}

Für die Projektdokumentation und die technische Ausarbeitung wurden verschiedene KI-gestützte Arbeitsschritte durchgeführt. Dabei standen vor allem die Recherche, die Strukturierung von Ideen und die sprachliche Überarbeitung im Vordergrund.

Im Bereich der Softwareentwicklung wurden Bezeichnungen im Code überprüft und verbessert. Dabei ging es vor allem um verständliche Namen für Klassen, Methoden, Variablen und fachliche Begriffe. Dadurch sollten die Struktur des Codes und die Lesbarkeit für alle Projektmitglieder verbessert werden.

Bei der Planung der Kundenkommunikation wurden mögliche Kommunikationswege wie SMS oder WhatsApp miteinander verglichen. Dabei wurde betrachtet, wie Kunden über geplante Liefertermine informiert werden können, wie eine Rückmeldung zum Termin erfolgen könnte und welche technischen Anforderungen sich daraus für das System ergeben.

Für das Thema Live-Tracking wurden grundlegende Umsetzungsmöglichkeiten untersucht. Dazu gehörten Überlegungen dazu, welche Standortinformationen für eine Lieferung relevant sind, wie diese Daten an das Backend übertragen werden können und wie eine verständliche Darstellung im Frontend aussehen könnte.

Auch einzelne Abschnitte der Dokumentation wurden sprachlich überarbeitet. Dabei wurden Formulierungen verständlicher gemacht, Inhalte besser gegliedert und technische Beschreibungen an den Stil einer Projektdokumentation angepasst.

Alle Ergebnisse wurden anschließend von der Projektgruppe geprüft, angepasst und in den jeweiligen Arbeitsbereich eingeordnet. Die fachlichen Entscheidungen, die technische Umsetzung und die finale Ausarbeitung erfolgten durch die Projektmitglieder.